% Chapter 4, Figure 9(b): percent error in y_b, F7 engine (scan: figures/ch4/fig09b.png).
% Curves: fig09.py, DIGITIZED from the 1973 art (the book does not give the F7 thrust curve or propellant mass, so
% the interval method (83)-(87) against eqs. (20), (21), (27), (28), (67) cannot be rerun); calibration
% fig09b.calib.json. Template: fig06a.tex.
\documentclass[11pt]{standalone}
\usepackage{tamrfig}
\begin{document}
\begin{tikzpicture}
  \begin{axis}[tamr,
      width=4.0in, height=2.5in,
      xmin=0.10, xmax=0.30, xtick={0.10,0.15,0.20,0.25,0.30}, minor xtick={0.125,0.175,0.225,0.275},
      xticklabel style={/pgf/number format/.cd, fixed, fixed zerofill, precision=2},
      ymin=-15, ymax=15, ytick distance=5,
      xlabel={$m_o$ (kg)}, ylabel={Percent error in $y_b$},
      legend pos=north east, legend style={nodes={inner sep=1pt}}]
    \draw[muted, line width=0.4pt] (axis cs:0.10,0) -- (axis cs:0.30,0);
    \addplot[series1] table[x=m0, y=fm_kmax, col sep=comma] {figures/v2/ch4/fig09b.csv};
    \addlegendentry{Fehskens-Malewicki}
    \addplot[series2] table[x=m0, y=cb_kmax, col sep=comma] {figures/v2/ch4/fig09b.csv};
    \addlegendentry{Caporaso-Bengen}
    \addplot[series1] table[x=m0, y=fm_kmin, col sep=comma] {figures/v2/ch4/fig09b.csv};
    \addplot[series2] table[x=m0, y=cb_kmin, col sep=comma] {figures/v2/ch4/fig09b.csv};
    % k labels with a leader to each method's curve for that k (end points on the curves: fig09b.csv)
    \node (kmax) at (axis cs:0.254,2.6) {\footnotesize $k_{\max}$};
    \draw[leader] (kmax.south west) -- (axis cs:0.2467,-1.976);
    \draw[leader] (kmax.south) -- (axis cs:0.2565,-6.604);
    \node (kmin) at (axis cs:0.145,-4.8) {\footnotesize $k_{\min}$};
    \draw[leader] (kmin.north east) -- (axis cs:0.1566,-1.159);
    \draw[leader] (kmin.south west) -- (axis cs:0.1397,-7.052);
  \end{axis}
\end{tikzpicture}
\end{document}
